\documentclass[12pt,a4paper]{article}
\usepackage{hojasemanal}

% ─── EDITAR AQUÍ ─────────────────────────────────────────────────────────────
\frase{¡Tú puedes!}  % opcional: frase al final de los ejercicios
% ─────────────────────────────────────────────────────────────────────────────

\begin{document}

    \encabezado

    \begin{indicaciones}
        \item Si no logras terminar una demostración, escribe hasta donde llegaste.
        \item Anota en la bitácora cuánto te tomó cada ejercicio.
        \item Ve apuntando en algún lugar las dudas que te surjan, para platicarlas en la próxima sesión.
    \end{indicaciones}

    \begin{mensaje}
        Hola Kaori, he aquí el formato creado única y exclusivamente de mí para ti. Vas a encontrar varias proposiciones que tendrás que demostrar con lo que vimos en la primera \enquote{sesión} (no sé cómo llamarla). Recuerda que tampoco es obligatorio, así que con calma, y por lo mismo trata de no usar IA. Y pues eso: recuerda que eres increíble, maravillosa y muy capaz, y sobre todo recuerda que te quiero mucho mucho.
    \end{mensaje}

    \parte{Demostración directa}

    \begin{ejercicio}[1]
        Demuestra que la suma de dos enteros pares es par.
    \end{ejercicio}

    \begin{ejercicio}[1]
        Demuestra que si $n$ es un entero impar, entonces $n^2$ es impar.
    \end{ejercicio}

    \begin{ejercicio}[1]
        Sean $a,b,c\in\Z$. Demuestra que si $a\mid b$ y $b\mid c$, entonces $a\mid c$.
    \end{ejercicio}

    \begin{ejercicio}[2]
        Sean $A$, $B$ y $C$ conjuntos con $A\subseteq B$. Demuestra que $A\cap C\subseteq B\cap C$.
    \end{ejercicio}

    \begin{ejercicio}[2]
        Demuestra que si $x$ es un número real con $x>0$, entonces
        \begin{equation*}
            x+\frac{1}{x}\geq2.
        \end{equation*}
    \end{ejercicio}

    \parte{Demostración por contrapositiva}

    \begin{ejercicio}[1]
        Demuestra que si $n$ es un entero y $n^2$ es par, entonces $n$ es par.
    \end{ejercicio}

    \begin{ejercicio}[2]
        Sean $a,b\in\Z$. Demuestra que si $ab$ es par, entonces $a$ es par o $b$ es par.
    \end{ejercicio}

    \begin{ejercicio}[2]
        Sean $a,b\in\Z$. Demuestra que si $a+b\geq15$, entonces $a\geq8$ o $b\geq8$.
    \end{ejercicio}

    \parte{Demostración por contradicción}

    \begin{ejercicio}[2]
        Demuestra que no existe $M\in\Z$ tal que $n\leq M$ para todo $n\in\Z$.
    \end{ejercicio}

    \begin{ejercicio}[2]
        Demuestra que $\sqrt{2}$ es irracional.
    \end{ejercicio}

    \begin{ejercicio}[2]
        Sean $a\in\Q$ y $b\in\R\setminus\Q$. Demuestra que $a+b$ es irracional.
    \end{ejercicio}

    \parte{Demostración por casos ($P\vee Q\implies R$)}

    \begin{ejercicio}[1]
        Sean $a,b\in\Z$. Demuestra que si $a$ es par o $b$ es par, entonces $ab$ es par.
    \end{ejercicio}

    \begin{ejercicio}[1]
        Sea $x\in\R$. Demuestra que si $x\leq-1$ o $x\geq1$, entonces $x^2\geq1$.
    \end{ejercicio}

    \begin{ejercicio}[2]
        Demuestra que $n^2+n$ es par para todo $n\in\Z$.
    \end{ejercicio}

    \begin{ejercicio}[2]
        Demuestra que $-\abs{x}\leq x\leq\abs{x}$ para todo $x\in\R$.
    \end{ejercicio}

    \parte{Reto (opcional)}

    \begin{ejercicio}[3]
        Demuestra que no existen $a,b\in\Z$ tales que $a^2=4b+3$.
    \end{ejercicio}

    \begin{ejercicio}[3]
        Demuestra que hay una infinidad de números primos.
    \end{ejercicio}

    \bitacora

\end{document}
